\section{A Core Language for Funnelling}
\label{sec:core}
% Metavariables are the letters of the grammar (X, T, C, e, Delta, Gamma, ...), decorated only with primes and digit subscripts:
% X, X', X'', X0, X3.
% The other subscripts are the indices of a sequence: X_i is the element at position i of X_1...X_n, and Xs is X_1...X_n.


\begin{figure}[H]
\centering
\[
\begin{array}{rrcl}
\textrm{Expressions} & e & ::= & x \mid \call{e}{\pl{T}}{m}{\pl{e}} \mid \explicitcall{e}{T}{\pl{T}}{m}{\pl{e}}
    \mid e\fstop f
    \mid \obj{C}{\pl{T}}{\pl{e}}\\
    &&&\mid \innerobj{C}{\pl{F}}{C'}{\pl{T}'}{\pl{e}}{\pl{M}}\\
    &&& \mid \ifthen{e_0}{T}{x}{e_1}{e_2}\\
\textrm{Types} & T & ::= & \gen{C}{\pl{T}} \mid X\\
\textrm{Funnels} & F & ::= & G \mid \funnel{X}\\
\textrm{Ground Types} & G & ::= & \gen{C}{\pl{G}}\\
\\
\textrm{Type Env.} & \Gamma & ::= & T_1\;x_1,\ldots,T_n\;x_n\\
\textrm{Bounds} & \Delta & ::= & \ext{X_1}{T_1},\ldots,\ext{X_n}{T_n}\\
%// Need T, not F, to be able to do things like (Java) class Bar<X,Y extends X> {}
\textrm{Declaration} & \mathsf{Dec} & ::= & \subclass{C}{\Delta}{C'}{\pl{T}}{\pl{M}}{\pl{S}}\\
\textrm{(Field) Slots} & S & ::= & T\;f\semi\\
\textrm{Methods} & M & ::= & \method{\Delta}{T_0}{m}{T_1\;x_1,\ldots,T_n\;x_n}{\ret{e}}\\
\textrm{Class Decl.} & K & ::= & \mathsf{Dec} \mid \innerobj{C}{\pl{F}}{C'}{\pl{T}}{\pl{e}}{\pl{M}}
%Dec::= class C<Bs> extends C'<Ts>{ Ms Ss} //assume standard constructor
%S ::= T f; //field/slot
%M ::= Sig{ return e; }  //no such a thing as an abstract method
%Sig::= <Bs> T m(x1:T1,..xn:Tn)
%G ::= x1:T1..xn:Tn
%v::=new C<Gs>(vs) | new C<Gs> extends C'<Gs'>(vs){Ms} // could be extends C<Gs>(vs){Ms}
%ctx::= [] 
%   | [].<Gs>m(es) | v.<Gs>m(vs [] es)
%   | [].[G]<Gs>m(es) | v.[G]<Gs>m(vs [] es)|
%   | [].f | new C<Gs>(vs [] es) 
%   | new C<Gs> extends C'<Gs'>(vs [] es){Ms} 
%   | if ([] instanceof G x){ e1 } else { e2 }
%
\end{array}
\]
\caption{Syntax}
\label{fig:syntax}
\end{figure}

\noindent
A program is a class table \decls and a main expression $e$.
We require the following well-formedness conditions:
\begin{enumerate}
    \item[(A1)] Class names are unique across top-level and local classes.
    \item[(A2)] Field names are unique across all classes.
    \item[(A3)] In a class, top-level or local, all method names are disjoint.
    \item[(A4)] In a method, all the parameters $x_1,\ldots,x_n$ are disjoint.
    \item[(A5)] In a $\Delta$, all the $X$ are disjoint.
    \item[(A6)] In a list of funnels $\pl{F}$, all the funnelled type variables are disjoint.
    \item[(A7)] There is no shadowing of $X$ or $x$, except for \this. No parameter and no variable of an \kw{if} is \this.
    \item[(A8)] \decls contains $\subclass{\Object}{}{\Object}{}{}{}$, and the chain of superclasses of every other top-level class reaches \Object.
\end{enumerate}
Well-formedness requirements A1-A8 are common requirements, such as uniqueness of names.
To simplify scoping, no shadowing is allowed.
A funnel $\funnel{X}$ does not bind $X$: it uses the type variable $X$ in scope.
A method of a local class binds \this, as in Java; this is the only form of shadowing.
Requirement A6 makes the substitutions of Figure~\ref{fig:substitution} well defined when their domain is a list of funnels.
Requirement A8 makes inheritance acyclic, except for \Object, which is its own superclass.
Thus rule \textsc{sub-inherit} can be applied to $\gen{\Object}{}$ forever, and rules \textsc{simpleFieldSuper} and \textsc{searchSuper} reduce a field access or a call on an \Object\ to itself.
This is harmless: our rules are not meant to be read as algorithms, and a well-typed expression never accesses a field or calls a method of \Object.
\\[1ex]
Well-formedness requirements B1-B3 below are about local classes:
\begin{enumerate}
    \item[(B1)] The main expression contains no local classes.
    \item[(B2)] Local classes cannot be extended: the superclass of every class, top-level or local, is a top-level class.
    \item[(B3)] In a local class $\innerobj{C}{\pl{F}}{C'}{\pl{T}}{\pl{e}}{\pl{M}}$, every funnelled type variable is in scope.
\end{enumerate}
Requirement B1 makes every local class occur in a method body: the judgement $\declared{K}$ and the function $\mathsf{scope}$ of Figure~\ref{fig:declarations} only search method bodies.
Requirements B2 and B3 are enforced by the type system, but we call them out because they are explicit design choices.
Requirement B1 is also enforced by rule \textsc{new-local}: a local class $C$ in the main expression is not declared, so no type $\gen{C}{\pl{T}}$ is well kinded.
Rule \textsc{new-local-java} of Figure~\ref{fig:java} does not check the kinding of the type of the local class; under rule \textsc{new-local-java}, requirement B1 is needed.

\begin{figure}[H]
\centering
\begin{mathpar}
\fbox{\declared{K}}\\
\inferrule[decl-top]{
  \mathsf{Dec}\in\decls
}{
  \declared{\mathsf{Dec}}
}
\and
\inferrule[decl-in-class]{
  \declared{\subclass{C}{\_}{\_}{\_}{\pl{M}}{\_}}\\\\
  \method{\_}{\_}{m}{\_}{\ret{e}}\in\pl{M}\\\\
  \innerobj{C_0}{\pl{F}}{C'}{\pl{T}}{\pl{e}}{\pl{M}'}\textrm{ occurs in }e
}{
  \declared{\innerobj{C_0}{\pl{F}}{C'}{\pl{T}}{\pl{e}}{\pl{M}'}}
}
\and
\inferrule[decl-in-local]{
  \declared{\innerobj{C}{\_}{\_}{\_}{\_}{\pl{M}}}\\\\
  \method{\_}{\_}{m}{\_}{\ret{e}}\in\pl{M}\\\\
  \innerobj{C_0}{\pl{F}}{C'}{\pl{T}}{\pl{e}}{\pl{M}'}\textrm{ occurs in }e
}{
  \declared{\innerobj{C_0}{\pl{F}}{C'}{\pl{T}}{\pl{e}}{\pl{M}'}}
}
\\
\fbox{$\mathsf{scope}(\gen{C}{\pl{T}})=\Delta;\pl{T}'$}
\and
\fbox{$\mathsf{scope}(\Delta;\pl{F};\pl{T})=\Delta';\pl{T}'$}
\end{mathpar}
\[
\begin{array}{rcl}
\mathsf{scope}(\gen{C}{\pl{T}}) & = & \Delta;\pl{T}\\
\multicolumn{3}{l}{\quad\textrm{where }\subclass{C}{\Delta}{\_}{\_}{\_}{\_}\in\decls}\\
\mathsf{scope}(\gen{C_0}{\pl{T}}) & = & \mathsf{scope}(\Delta',\Delta'';\pl{F};\pl{T})\\
\multicolumn{3}{l}{\quad\textrm{where }\declared{\subclass{C}{\_}{\_}{\_}{\pl{M}}{\_}}\textrm{ or }\declared{\innerobj{C}{\_}{\_}{\_}{\_}{\pl{M}}}}\\
\multicolumn{3}{l}{\quad\phantom{\textrm{where }}\method{\Delta''}{\_}{m}{\_}{\ret{e}}\in\pl{M}}\\
\multicolumn{3}{l}{\quad\phantom{\textrm{where }}\innerobj{C_0}{\pl{F}}{\_}{\_}{\_}{\_}\textrm{ occurs in }e}\\
\multicolumn{3}{l}{\quad\phantom{\textrm{where }}\mathsf{scope}(\generic{C})=\Delta';\_}
\end{array}
\]
\[
\begin{array}{rcl}
\mathsf{scope}(\Delta;\emptyset;\emptyset) & = & \emptyset;\emptyset\\
\mathsf{scope}(\Delta;\funnel{X},\pl{F};T,\pl{T}) & = & (\ext{X}{\Delta(X)};T)\cdot\mathsf{scope}(\Delta;\pl{F};\pl{T})\\
\mathsf{scope}(\Delta;G,\pl{F};G,\pl{T}) & = & \mathsf{scope}(\Delta;\pl{F};\pl{T})\\[1ex]
(\Delta;\pl{T})\cdot(\Delta';\pl{T}') & = & \Delta,\Delta';\pl{T},\pl{T}'
\end{array}
\]
\begin{mathpar}
\fbox{$\generic{C}=T$}
\and
\fbox{$|F|=T$}
\end{mathpar}
\[
\begin{array}{rcl}
\generic{C} & = & \gen{C}{\mathrm{dom}(\Delta)}\\
\multicolumn{3}{l}{\quad\textrm{where }\subclass{C}{\Delta}{\_}{\_}{\_}{\_}\in\decls}\\
\generic{C} & = & \gen{C}{|\pl{F}|}\\
\multicolumn{3}{l}{\quad\textrm{where }\declared{\innerobj{C}{\pl{F}}{\_}{\_}{\_}{\_}}}\\[1ex]
|\funnel{X}| & = & X\\
|G| & = & G
\end{array}
\]
\caption{Declarations}
\label{fig:declarations}
\end{figure}

By well-formedness requirement A1, a class name $C$ identifies at most one $K$ with $\declared{K}$.
For a class type $\gen{C}{\pl{T}}$, $\mathsf{scope}(\gen{C}{\pl{T}})=\Delta_0;\pl{T}''$ gives the bounds $\Delta_0$ of the type variables that the methods of $C$ see, without the type variables of the methods themselves.
It also gives the types $\pl{T}''$ that the arguments $\pl{T}$ assign to the type variables in $\Delta_0$.
For a top-level class, $\Delta_0$ is the bounds of the class and $\pl{T}''=\pl{T}$.
For a top-level class, $\mathsf{scope}(\gen{C}{\pl{T}})$ does not check that $\pl{T}$ has as many types as $\Delta_0$ has type variables.
Rule \textsc{wf-class} checks it: its premise $\subtype{\Delta}{\pl{T}''}{\tsubst{\mathrm{cod}(\Delta_0)}{\Delta_0}{\pl{T}''}}$ is a judgement on sequences, which requires sequences of the same length.
For a local class $C_0$ with funnels $\pl{F}$, the bounds $\Delta',\Delta''$ are the bounds of the type variables in scope where $C_0$ is written.
For a local class, $\mathsf{scope}(\Delta;\pl{F};\pl{T})$ is undefined when $\pl{F}$ and $\pl{T}$ have different lengths.
\\[1ex]
For a type variable $X$, $\Delta(X)$ is the bound of $X$ in $\Delta$.
$\mathrm{dom}(\Delta)$ is the sequence of the type variables in $\Delta$, and $\mathrm{cod}(\Delta)$ is the sequence of the bounds in $\Delta$.
$\ftv{T}$ is the set of the type variables occurring in $T$; for a sequence, $\mathsf{FTV}$ is the union over its elements.
The figures leave the following conventions implicit.
A sequence $e_1,\ldots,e_n$, or $\pl{e}$, separates its elements with commas; a comma also concatenates two sequences, and either sequence may be empty.
A function or a judgement applied to sequences stands for one application at every position, and requires the sequences to have the same length:
$|\pl{F}|$ is $|F_1|,\ldots,|F_n|$;
$\kind{\Delta}{\pl{T}}$ is $\kind{\Delta}{T_i}$ for every $i$;
$\subtype{\Delta}{T_1,\ldots,T_n}{T'_1,\ldots,T'_n}$ is $\subtype{\Delta}{T_i}{T'_i}$ for every $i$;
$\typing{\Delta}{\Gamma}{e_1,\ldots,e_n}{T_1,\ldots,T_n}$ is $\typing{\Delta}{\Gamma}{e_i}{T_i}$ for every $i$;
$\mathsf{overrideOk}(T,\pl{M})$ is $\mathsf{overrideOk}(T,M)$ for every $M$ in $\pl{M}$.
Finally, $m\not\in\pl{M}$ holds when no method in $\pl{M}$ has name $m$.
\\[1ex]
The function $\mathsf{scope}(\Delta;\pl{F};\pl{T})$ goes through the funnels $\pl{F}$ and the arguments $\pl{T}$ together, one position at a time, where $\emptyset$ is the empty sequence.
At the position of a funnel $\funnel{X}$, it adds $\ext{X}{\Delta(X)}$ to the bounds and the argument to the types.
At the position of a ground funnel $G$, the argument must be $G$ itself, and nothing is added.
The operator $\cdot$ concatenates two results of $\mathsf{scope}$, component by component.
An expression $e'$ occurs in $e$ when $e=\mathcal{C}[e']$, for a context $\mathcal{C}$ of Figure~\ref{fig:contexts}.
Contexts $\mathcal{C}$ do not enter the methods of a local class; thus $\mathsf{scope}(\gen{C}{\pl{T}})$ has at most one value.
The reduction of Figure~\ref{fig:opsem} uses the values $v$ and the evaluation contexts $\mathcal{E}$ of Figure~\ref{fig:contexts}.
\\[1ex]
Rules \textsc{simpleFound} and \textsc{innerFound} use the method $\mathsf{fresh}(M)$:
the method $M$ where every type variable and every variable bound in $M$ is renamed into a fresh name, that is a name occurring nowhere in the program and in the reduced expression; distinct binders get distinct fresh names.
The type variables and variables bound in $M$ are the type variables of its signature, its parameters, and, inside its body, the variable of every \kw{if} and the type variables and parameters of every method of a local class.
The variable \this is never renamed.
Renaming is needed because a value of a local class carries the code of its methods.
For example, consider a local class \QQ{L} whose method \QQ{m} has a parameter \QQ{x} and calls \QQ{m} on a captured value of class \QQ{L}.
The body of a call of \QQ{m} then contains the captured value, and the captured value contains a copy of the method \QQ{m} binding \QQ{x} again.
Without renaming, rule \textsc{innerFound} would replace \QQ{x} also in the method \QQ{m} of the captured value.
With renaming, the domain of every substitution is disjoint from the names bound inside the substituted expression.
Only values of local classes carry code, by requirement B1 local classes occur only in method bodies,
and a method body enters the reduced expression only through rules \textsc{simpleFound} and \textsc{innerFound}, renamed by $\mathsf{fresh}$.
A substitution can copy a value, so several copies of a value can bind the same names.
Those names are never in the domain of a later substitution.
Rules \textsc{simpleFound} and \textsc{innerFound} substitute only names that $\mathsf{fresh}$ has just created.
Rule \textsc{ifThen} substitutes the variable of an \kw{if}; this variable is either renamed by $\mathsf{fresh}$ or written in the main expression, and no value binds it.
Outside of values, the names bound inside the reduced expression are not shadowed, by requirement A7.
Thus a substitution never reaches a type variable or a variable bound inside the substituted expression.

\begin{figure}[H]\small
    \centering
\begin{mathpar}
\fbox{$\step{e}{e'}$}\\
\inferrule[simpleFieldFound]{
    \subclass{C}{\_}{\_}{\_}{\_}{T_1\;f_1\semi\ldots T_n\;f_n\semi}\in\decls\\
    f_i\in f_1,\ldots,f_n
}{
    \step{
        \obj{C}{\pl{G}}{v_1,\ldots,v_n,\pl{v}}\fstop f_i
    }{
        v_i
    }
}
\and
\inferrule[simpleFieldSuper]{
    \subclass{C}{\Delta}{C'}{\pl{T}}{\_}{T_1\;f_1\semi\ldots T_n\;f_n\semi}\in\decls\\
    f\not\in f_1,\ldots,f_n
}{
    \step{
        \obj{C}{\pl{G}}{v_1,\ldots,v_n,\pl{v}}\fstop f
    }{
        \obj{C'}{\tsubst{\pl{T}}{\Delta}{\pl{G}}}{\pl{v}}\fstop f
    }
}
\and
%Note: ok because inner classes currently don't have their own fields
\inferrule[innerFieldSuper]{
}{
    \step{
        \innerobj{C}{\pl{G}}{C'}{\pl{G}'}{\pl{v}}{\pl{M}}\fstop f
    }{
        \obj{C'}{\pl{G}'}{\pl{v}}\fstop f
    }
}
\and
\inferrule[ifThen]{
    v\in\{\obj{C}{\pl{G}}{\_}, \innerobj{C}{\pl{G}}{\_}{\_}{\_}{\_}\}\\
    \emptyset\vdash \gen{C}{\pl{G}} <: \gen{C'}{\_}
    %\match{v}{C\langle\_\rangle} = \mathsf{true}
}{
    \step{\ifthen{v}{\gen{C'}{\pl{T}}}{x}{e_1}{e_2}}
    {\tsubst{e_1}{x}{v}}
}
\and
\inferrule[ifElse]{
    %v\not\in\{\obj{C}{\_}{\_}, \innerobj{C}{\_}{\_}{\_}{\_}{\_}\}
    %\match{v}{C\langle\_\rangle} = \mathsf{false}
    v\in\{\obj{C}{\pl{G}}{\_}, \innerobj{C}{\pl{G}}{\_}{\_}{\_}{\_}\}\\
    \emptyset\not\vdash \gen{C}{\pl{G}} <: \gen{C'}{\_}
}{
    \step{\ifthen{v}{\gen{C'}{\pl{T}}}{x}{e_1}{e_2}}
    {e_2}
}
% The formal model has only the erased type test of Java: rules ifThen and ifElse compare class names only.
% Reified type tests are left out of the formal model on purpose, and are discussed informally in Section 2:
% under funnelled typing, an erased and a reified type test give the same answer on every well-typed program;
% under Java typing, only a deep reified type test, also checking the type arguments of the superclasses, catches the unsoundness.
% Neither Java nor funnelling relies on reification.
\and
\inferrule[simpleAddSuper]{
    v = \obj{C}{\pl{G}}{\_}
}{
    \step{
        \call{v}{\pl{G}'}{m}{\pl{v}} 
    }{
        \explicitcall{v}{\gen{C}{\pl{G}}}{\pl{G}'}{m}{\pl{v}}
    }
}
\and
\inferrule[simpleFound]{
    \subclass{C}{\Delta}{\_}{\_}{\pl{M}}{\_}\in\decls\\
    M\in\pl{M}\\
    \mathsf{fresh}(M)=\method{\Delta'}{T_0}{m}{T_1\;x_1,\ldots,T_n\;x_n}{\ret{e_0}}\\
    e = \tsubst{\tsubst{e_0}{\Delta,\Delta'}{\pl{G},\pl{G}'}}{\this,x_1,\ldots,x_n}{v_0,\ldots,v_n}
}{
    \step{\explicitcall{v_0}{\gen{C}{\pl{G}}}{\pl{G}'}{m}{v_1,\ldots,v_n}}{e}
}
%% ^^ Note: for local classes, this syntax skips the value entirely (this rule is not applicable if C<Gs> is meth local
\and
\inferrule[innerFound]{
    v_0 = \innerobj{\_}{\_}{\_}{\_}{\_}{\pl{M}}\\
    M\in\pl{M}\\
    \mathsf{fresh}(M)=\method{\Delta'}{T_0}{m}{T_1\;x_1,\ldots,T_n\;x_n}{\ret{e_0}}\\
    e = \tsubst{\tsubst{e_0}{\Delta'}{\pl{G}'}}{\this,x_1,\ldots,x_n}{v_0,\ldots,v_n}
}{
    \step{\call{v_0}{\pl{G}'}{m}{v_1,\ldots,v_n}}{e}
}
\and
\inferrule[innerSuper]{
    v_0 = \innerobj{\_}{\_}{C}{\pl{G}}{\_}{\pl{M}}\\
    m\not\in \pl{M}
}{
    \step{\call{v_0}{\pl{G}'}{m}{v_1,\ldots,v_n}}{
        \explicitcall{v_0}{\gen{C}{\pl{G}}}{\pl{G}'}{m}{v_1,\ldots,v_n}
    }
}
\and
\inferrule[searchSuper]{
    \subclass{C}{\Delta}{C'}{\pl{T}}{\pl{M}}{\_}\in\decls\\
    m\not\in \pl{M}
}{
    \step{
        \explicitcall{v}{\gen{C}{\pl{G}}}{\pl{G}'}{m}{\pl{v}}
    }{
        \explicitcall{v}{\gen{C'}{\tsubst{\pl{T}}{\Delta}{\pl{G}}}}{\pl{G}'}{m}{\pl{v}}
    }
}
\and
\inferrule[ctx]{
    \step{e}{e'}
}{
    \step{\ctx{e}}{\ctx{e'}}
}
\end{mathpar}
    \caption{Operational Semantics}
    \label{fig:opsem}
\end{figure}

%Not totally sure about Ts vs Gs vs Fs in LHS of match
%\[
%    \begin{array}{rcl}
%    \match{\obj{C}{Gs}{vs}}{C\langle Gs\rangle} & = & \mathsf{true}\\
%    \match{\obj{C}{Gs}{vs}}{C\langle Gs'\rangle} & = & \mathsf{false}\\
%    \match{\obj{C}{Gs}{v_1\ldots v_n\;vs}}{D\langle Gs'\rangle} & = & \match{\obj{C'}{\tsubst{Ts}{Bs}{Gs}}{vs}}{D\langle Gs'\rangle} \quad \textrm{if}~\subclass{C}{Bs}{C'}{Ts}{Ms}{fs}\in\decls\\
%    \match{\obj{C}{Gs}{v_1\ldots v_n\;vs}}{D\langle Gs'\rangle} & = & \mathsf{false} \quad \textrm{otherwise}\\
%    \match{\innerobj{C}{Gs}{C'}{Gs'}{vs}{Ms}}{C\langle Gs\rangle} &=& \mathsf{true}\\
%    \match{\innerobj{C}{Gs}{C'}{Gs'}{vs}{Ms}}{D\langle Gs''\rangle} &=& \match{\obj{C'}{Gs'}{vs}}{D\langle Gs''\rangle}\\
%    \hline
%    \match{\obj{C}{Gs}{vs}}{C\langle\_\rangle} & = & \mathsf{true}\\
%    \match{\obj{C}{Gs}{v_1\ldots v_n\;vs}}{D\langle\_\rangle} & = & \match{\obj{C'}{\tsubst{Ts}{Bs}{Gs}}{vs}}{D\langle\_\rangle} \quad \textrm{if}~\subclass{C}{Bs}{C'}{Ts}{Ms}{fs}\in\decls\\
%    \match{\obj{C}{Gs}{v_1\ldots v_n\;vs}}{D\langle\_\rangle} & = & \mathsf{false} \quad \textrm{otherwise}\\
%    \match{\innerobj{C}{Gs}{C'}{Gs'}{vs}{Ms}}{C\langle\_\rangle} &=& \mathsf{true}\\
%    \match{\innerobj{C}{Gs}{C'}{Gs'}{vs}{Ms}}{D\langle\_\rangle} &=& \match{\obj{C'}{Gs'}{vs}}{D\langle\_\rangle}\\
%    \end{array}
%\]

\begin{figure}[H]\small
    \centering
    \begin{mathpar}
\fbox{\typing{\Delta}{\Gamma}{e}{T}}\\
\inferrule[subs]{
	\typing{\Delta}{\Gamma}{e}{T_1}\\
	\subtype{\Delta}{T_1}{T_2}\\
	\kind{\Delta}{T_2}
}{
	\typing{\Delta}{\Gamma}{e}{T_2}
}
\and
\inferrule[var]{ }{\typing{\Delta}{\Gamma}{x}{\Gamma(x)}}
\and
\inferrule[new-object]{ }{
  \typing{\Delta}{\Gamma}{\obj{\Object}{}{}}{\gen{\Object}{}}
}
\and
\inferrule[declField]{
  \typing{\Delta}{\Gamma}{e}{\classtype{C}{\pl{T}}}\\\\
  \subclass{C}{\Delta'}{\_}{\_}{\_}{\pl{S}}\in\decls\\\\
  T\;f\semi\in\pl{S}
}{
  % thanks to subsumption, no need of multiple rules (C<Ts> can simply be set to supertype of relevance, via subsumption)
  \typing{\Delta}{\Gamma}{e\fstop f}{\tsubst{T}{\Delta'}{\pl{T}}}
  %\subclass{C_0}{\Delta_0}{C'}{\overline{T'}}{\overline{M}}{\overline{S}}\in\decls
}
\and
\inferrule[call-top]{
  \typing{\Delta}{\Gamma}{e_0}{\classtype{C_0}{\pl{T}'}}\\
  \kind{\Delta}{\pl{T}''}\\\\
  \declared{\subclass{C_0}{\Delta_0}{\_}{\_}{\pl{M}}{\_}}\\\\
  \method{\Delta_1}{T_0}{m}{T_1\;x_1,\ldots,T_n\;x_n}{\_} \in \pl{M}\\\\
  \subtype{\Delta}{\pl{T}''}{\tsubst{\mathrm{cod}(\Delta_1)}{\Delta_0,\Delta_1}{\pl{T}',\pl{T}''}}\\\\
  \typing{\Delta}{\Gamma}{e_1,\ldots,e_n}{\tsubst{T_1,\ldots,T_n}{\Delta_0,\Delta_1}{\pl{T}',\pl{T}''}}
}{
  \typing{\Delta}{\Gamma}{\call{e_0}{\pl{T}''}{m}{e_1,\ldots,e_n}}{\tsubst{T_0}{\Delta_0,\Delta_1}{\pl{T}',\pl{T}''}}
}
\and
\inferrule[call-local]{
  \typing{\Delta}{\Gamma}{e_0}{\classtype{C_0}{\pl{T}'}}\\
  \kind{\Delta}{\pl{T}''}\\\\
  \declared{\innerobj{C_0}{\pl{F}}{\_}{\_}{\_}{\pl{M}}}\\\\
  \method{\Delta_1}{T_0}{m}{T_1\;x_1,\ldots,T_n\;x_n}{\_} \in \pl{M}\\\\
  \subtype{\Delta}{\pl{T}''}{\tsubst{\mathrm{cod}(\Delta_1)}{\pl{F},\Delta_1}{\pl{T}',\pl{T}''}}\\\\
  \typing{\Delta}{\Gamma}{e_1,\ldots,e_n}{\tsubst{T_1,\ldots,T_n}{\pl{F},\Delta_1}{\pl{T}',\pl{T}''}}
}{
  \typing{\Delta}{\Gamma}{
    \call{e_0}{\pl{T}''}{m}{e_1,\ldots,e_n}
  }{\tsubst{T_0}{\pl{F},\Delta_1}{\pl{T}',\pl{T}''}}
}
\and
\inferrule[call-lb-runtime]{
  \typing{\Delta}{\Gamma}{e_0}{\gen{C}{\pl{T}'''}}\\
  \kind{\Delta}{\gen{C}{\pl{T}'''}}\\
  \kind{\Delta}{\pl{T}''}\\\\
  \subclass{C}{\_}{\_}{\_}{\_}{\_}\in\decls\\
  \subtype{\Delta}{\gen{C}{\pl{T}'''}}{\gen{C_0}{\pl{T}'}}\\\\
  \declared{\subclass{C_0}{\Delta_0}{\_}{\_}{\pl{M}}{\_}}\\
  \method{\Delta_1}{T_0}{m}{T_1\;x_1,\ldots,T_n\;x_n}{\_} \in \pl{M}\\\\
  \subtype{\Delta}{\pl{T}''}{\tsubst{\mathrm{cod}(\Delta_1)}{\Delta_0,\Delta_1}{\pl{T}',\pl{T}''}}\\\\
  \typing{\Delta}{\Gamma}{e_1,\ldots,e_n}{\tsubst{T_1,\ldots,T_n}{\Delta_0,\Delta_1}{\pl{T}',\pl{T}''}}
}{
  \typing{\Delta}{\Gamma}{
    \explicitcall{e_0}{\gen{C}{\pl{T}'''}}{\pl{T}''}{m}{e_1,\ldots,e_n}
  }{\tsubst{T_0}{\Delta_0,\Delta_1}{\pl{T}',\pl{T}''}}
}
\and
\inferrule[new-top]{
  \declared{\subclass{C}{\Delta'}{C'}{\pl{T}'}{\_}{T_1\;f_1\semi\ldots T_n\;f_n\semi}}\\\\
  \typing{\Delta}{\Gamma}{\obj{C'}{\tsubst{\pl{T}'}{\Delta'}{\pl{T}}}{\pl{e}}}{\_}\\
  \typing{\Delta}{\Gamma}{e_1,\ldots,e_n}{\tsubst{T_1,\ldots,T_n}{\Delta'}{\pl{T}}}\\
  \kind{\Delta}{\gen{C}{\pl{T}}}
}{
  \typing{\Delta}{\Gamma}{\obj{C}{\pl{T}}{e_1,\ldots,e_n,\pl{e}}}{\gen{C}{\pl{T}}}
}
\and
\inferrule[new-local]{
  \mathsf{scope}(\Delta;\pl{F};|\pl{F}|)=\Delta';\_\\
  \ftv{\mathrm{cod}(\Delta')}\subseteq\mathrm{dom}(\Delta')\\
  \kind{\Delta'}{\gen{C}{|\pl{F}|}}\\\\
  \typing{\Delta}{\Gamma}{\obj{C'}{\pl{T}}{\pl{e}}}{\_}\\
  \kind{\Delta'}{\gen{C'}{\pl{T}}}\\\\
  \Gamma'=\{ T\;x \mid T\;x\in\Gamma \land x\neq\this \land \ftv{T}\subseteq\ftv{|\pl{F}|} \}\\\\ % remove all the x:T when FTV(T) not a subset of FTV(T1..Tn)
  \mok{\Delta'}{\Gamma',\gen{C}{|\pl{F}|}\;\this}{\pl{M}}\\\\
  \forall C'',\pl{T}'\ldotp\subtype{\Delta'}{\gen{C'}{\pl{T}}}{\gen{C''}{\pl{T}'}}\Rightarrow
    \mathsf{overrideOk}(\gen{C''}{\pl{T}'},\pl{M})
}{
  \typing{\Delta}{\Gamma}{
    \innerobj{C}{\pl{F}}{C'}{\pl{T}}{\pl{e}}{\pl{M}}
  }{\gen{C}{|\pl{F}|}}
}
%
%  Bs;G |- new C'<Ts'>(es) : _//incredibly subtle: is Bs;G or Bs';G' // need Bs;G for es that internally use non-funnelled types but reduce to terms with funnelled types
%  Bs' |- C'<Ts'> : Ok // Not redundant with first line, b/c need to ensure supertype only depends on funnelled things
%  Ti= Xi if Fi = Xi funnel //those two lines do declare T1..Tn
%  Ti= Gi if Fi = Gi        //those two lines do declare T1..Tn
%  Bs' |- C<T1..Tn> : Ok
%  Bs'= {X extends T | X extends T in Bs, Fi=X funnel} // Bs' = Bs \ dom(F1..Fn)
%  forall Fi = Xi funnel, Xi extends _ in Bs
%  G' = G/FTV(T1..Tn) //remove all the x:T when FTV(T) not a subset of FTV(T1..Tn)
%  Bs';G',this:C<T1..Tn> |- M1 : Ok
%  ..
%  Bs';G',this:C<T1..Tn> |- Mk : Ok
%-------------------------------------------------------------------------
%  Bs;G,this:_ |- new C<F1..Fn> extends C'<Ts'>(es){M1..Mk} : C<T1..Tn>
%
%Note: we still need formal definition of substitution e[Xs=Ts]; it is crucial to go in the details for e of form meth local with funnelling
%We will replace every occurrence of X when doing substitution, *including* in method signatures in a local class
%Should be a clearer when substitution is formalized
%
\and
\inferrule[if-instanceof]{
  \typing{\Delta}{\Gamma}{e_0}{T_0}\\
  \kind{\Delta}{T}\\
  \cast{\Delta}{T_0}{T}\\\\
  \typing{\Delta}{\Gamma,T\;x}{e_1}{T'}\\
  \typing{\Delta}{\Gamma}{e_2}{T'}
}{
  \typing{\Delta}{\Gamma}{\ifthen{e_0}{T}{x}{e_1}{e_2}}{T'}
}
    \end{mathpar}
    \caption{Typing}
    \label{fig:typing}
\end{figure}

\begin{figure}[H]\small
    \centering
    \begin{mathpar}
\fbox{$\vdash\decls\;e~\mathrm{Ok}$}\\
\inferrule[prog-ok]{
  \forall\mathsf{Dec}\in\decls\ldotp\mathsf{Dec}~\mathrm{Ok}\\
  \typing{\emptyset}{\emptyset}{e}{\_}
}{
  \vdash\decls\;e~\mathrm{Ok}
}
\\
\fbox{$\mathsf{Dec}~\mathrm{Ok}$}\\
\inferrule[class-ok]{
  \kind{\Delta}{\mathrm{cod}(\Delta)}\\
  \forall T\;f\semi\in\pl{S}\ldotp\kind{\Delta}{T}\\\\
  \subclass{C'}{\_}{\_}{\_}{\_}{\_}\in\decls\\
  \kind{\Delta}{\gen{C'}{\pl{T}}}\\\\
  \mok{\Delta}{\gen{C}{\mathrm{dom}(\Delta)}\;\this}{\pl{M}}\\\\
  \forall C'',\pl{T}'\ldotp\subtype{\Delta}{\gen{C'}{\pl{T}}}{\gen{C''}{\pl{T}'}}\Rightarrow
    \mathsf{overrideOk}(\gen{C''}{\pl{T}'},\pl{M})
}{
  \subclass{C}{\Delta}{C'}{\pl{T}}{\pl{M}}{\pl{S}}~\mathrm{Ok}
}
\\
\fbox{\mok{\Delta}{\Gamma}{M}}\\
\inferrule[method-ok]{
  \kind{\Delta,\Delta'}{\mathrm{cod}(\Delta')}\\
  \kind{\Delta,\Delta'}{T_0,\ldots,T_n}\\\\
  \typing{\Delta,\Delta'}{\Gamma,T_1\;x_1,\ldots,T_n\;x_n}{e}{T_0}
}{
  \mok{\Delta}{\Gamma}{\method{\Delta'}{T_0}{m}{T_1~x_1,\ldots,T_n~x_n}{\ret{e}}}
}
%
%  Bs,Bs';G,x1:T1,..xn:Tn |- e : T0
%  Bs,Bs'|- T0 : Ok .. Bs|- Tn : Ok
%------------------------------------------------------------------------
%  Bs;G |- <Bs'> T0 m(T1 x1..Tn xn){return e;} : Ok
\\
\fbox{$\mathsf{overrideOk}(T,M)$}\\
\inferrule[override-new]{
  \subclass{C}{\_}{\_}{\_}{\pl{M}}{\_}\in\decls\\
  m\not\in\pl{M}
}{
  \mathsf{overrideOk}(\gen{C}{\pl{T}},\method{\_}{\_}{m}{\_}{\_})
}
\and
\inferrule[override-same]{
  \subclass{C}{\Delta}{\_}{\_}{\pl{M}}{\_}\in\decls\\\\
  \method{\Delta''}{T''_0}{m}{T''_1\;x'_1,\ldots,T''_n\;x'_n}{\_}\in\pl{M}\\\\
  \tsubst{\mathrm{cod}(\Delta'')}{\Delta,\Delta''}{\pl{T},\mathrm{dom}(\Delta')}=\mathrm{cod}(\Delta')\\\\
  \tsubst{T''_0,\ldots,T''_n}{\Delta,\Delta''}{\pl{T},\mathrm{dom}(\Delta')}=T'_0,\ldots,T'_n
}{
  \mathsf{overrideOk}(\gen{C}{\pl{T}},\method{\Delta'}{T'_0}{m}{T'_1\;x_1,\ldots,T'_n\;x_n}{\_})
}
    \end{mathpar}
    \caption{Typing of programs, classes and methods}
    \label{fig:classtyping}
\end{figure}

Rule \textsc{class-ok} types the methods of a top-level class with \this of type $\gen{C}{\mathrm{dom}(\Delta)}$.
A method that a class declares again must keep the signature it has in every superclass, up to the names of its own type variables:
\textsc{call-top} reads the signature from the class chosen by subsumption, while the reduction runs the method of the class of the receiver.

\begin{figure}[H]\small
    \centering
    \begin{mathpar}
\fbox{\kind{\Delta}{T}}\\
\inferrule[wf-var]{
  \ext{X}{\_}\in\Delta
}{
  \kind{\Delta}{X}
}
\and
\inferrule[wf-class]{
  \mathsf{scope}(\gen{C}{\pl{T}})=\Delta_0;\pl{T}''\\
  \kind{\Delta}{\pl{T}}\\\\
  \subtype{\Delta}{\pl{T}''}{\tsubst{\mathrm{cod}(\Delta_0)}{\Delta_0}{\pl{T}''}}
}{
  \kind{\Delta}{\gen{C}{\pl{T}}}
}
\\
\fbox{\subtype{\Delta}{T}{T'}}\\
\inferrule[sub-refl]{ }{\subtype{\Delta}{T}{T}}
\and
\inferrule[sub-inherit]{
  \subclass{C}{\Delta_0}{C'}{\pl{T}'}{\_}{\_}\in\decls\\\\
  \subtype{\Delta}{\tsubst{\gen{C'}{\pl{T}'}}{\Delta_0}{\pl{T}}}{T''}
}{
  \subtype{\Delta}{\gen{C}{\pl{T}}}{T''}
}
\and
\inferrule[sub-inherit-local]{
  \declared{\innerobj{C}{\pl{F}}{C'}{\pl{T}'}{\_}{\_}}\\\\
  \subtype{\Delta}{\tsubst{\gen{C'}{\pl{T}'}}{\pl{F}}{\pl{T}}}{T''}
}{
  \subtype{\Delta}{\gen{C}{\pl{T}}}{T''}
}
\and
\inferrule[sub-var]{
  \ext{X}{T}\in\Delta\\
  \subtype{\Delta}{T}{T'}
}{
  \subtype{\Delta}{X}{T'}
}
    \end{mathpar}
    \caption{Kinding and Subtyping}
    \label{fig:kinding}
\end{figure}

Rule \textsc{wf-class} kinds both top-level and local classes, including a local class written in a method of another local class.
Every type in $\pl{T}''$ is checked against the bound of its type variable in $\Delta_0$.
A funnelled type variable that is not in scope where its local class $C$ is written leaves $\mathsf{scope}(\gen{C}{\pl{T}})$ undefined; thus no type $\gen{C}{\pl{T}}$ is well kinded.

\begin{figure}[H]\small
    \centering
    \begin{mathpar}
\fbox{\cast{\Delta}{T}{T'}}\\
\inferrule[cast-up]{
  \subtype{\Delta}{T}{T'}
}{
  \cast{\Delta}{T}{T'}
}
\and
\inferrule[cast-var]{
  \ext{X}{T}\in\Delta\\
  \cast{\Delta}{T}{T'}
}{
  \cast{\Delta}{X}{T'}
}
\and
\inferrule[cast-down]{
  \subtype{\Delta}{\gen{C}{\pl{T}}}{\gen{C'}{\pl{T}_0}}\\
  \generic{C}=\gen{C}{\pl{T}''}\\\\
  \subtype{\emptyset}{\gen{C}{\pl{T}''}}{\gen{C'}{\pl{T}'}}\\
  \ftv{\pl{T}''}\subseteq\ftv{\pl{T}'}
}{
  \cast{\Delta}{\gen{C'}{\pl{T}_0}}{\gen{C}{\pl{T}}}
}
    \end{mathpar}
    \caption{Checked Casts}
    \label{fig:cast}
\end{figure}

\begin{figure}[H]\footnotesize
    \centering
    \setlength{\arraycolsep}{3pt}
\[
\newcommand{\s}[1]{\tsubst{#1}{\pl{X}}{\pl{T}}}
\begin{array}{rcl}
\s{x} & = & x\\[1ex]
\s{\call{e}{\pl{T}'}{m}{\pl{e}'}} & = & \call{\s{e}}{\s{\pl{T}'}}{m}{\s{\pl{e}'}}\\[1ex]
\s{\explicitcall{e}{T'}{\pl{T}''}{m}{\pl{e}'}} & = &
  \twolines{\s{e}\fstop\OS\s{T'}\CS\OA\s{\pl{T}''}\CA}{m\OR\s{\pl{e}'}\CR}\\[3ex]
\s{e\fstop f} & = & \s{e}\fstop f\\[1ex]
\s{\obj{C}{\pl{T}'}{\pl{e}}} & = & \obj{C}{\s{\pl{T}'}}{\s{\pl{e}}}\\[1ex]
\twolines{\kw{if}\;\OR e_0\;\kw{instanceof}\;T'\;x\CR}{\OC\;e_1\;\CC\;\kw{else}\;\OC\;e_2\;\CC\;\s{}} & = &
  \twolines{\kw{if}\;\OR\s{e_0}\;\kw{instanceof}\;\s{T'}\;x\CR}{\OC\;\s{e_1}\;\CC\;\kw{else}\;\OC\;\s{e_2}\;\CC}\\[3ex]
\twolines{\kw{new}\;\gen{C}{\pl{F}}\;\kw{extends}\;\gen{C'}{\pl{T}'}}{\OR\pl{e}\CR\OC\pl{M}\CC\;\s{}} & = &
  \twolines{\kw{new}\;\gen{C}{\s{\pl{F}}}\;\kw{extends}\;\gen{C'}{\s{\pl{T}'}}}{\OR\s{\pl{e}}\CR\OC\s{\pl{M}}\CC}\\[3ex]
\s{\funnel{X'}} & = & \funnel{X'} \quad\textrm{if }X'\not\in\pl{X}\\
\tsubst{\funnel{X_i}}{X_1,\ldots,X_n}{T_1,\ldots,T_n} & = & T_i\\
\s{G} & = & G\\[1ex]
\twolines{\methodsig{\Delta}{T'_0}{m}{T'_1\;x_1,\ldots,T'_k\;x_k}}{\OC\;\ret{e}\;\CC\;\s{}} & = &
  \twolines{\methodsig{\s{\Delta}}{\s{T'_0}}{m}{\s{T'_1}\;x_1,\ldots,\s{T'_k}\;x_k}}{\OC\;\ret{\s{e}}\;\CC}\\[3ex]
\s{(\ext{X'}{T'})} & = & \ext{X'}{\s{T'}}\\[1ex]
\s{\gen{C}{\pl{T}'}} & = & \gen{C}{\s{\pl{T}'}}\\
\s{X'} & = & X' \quad\textrm{if }X'\not\in\pl{X}\\
\tsubst{X_i}{X_1,\ldots,X_n}{T_1,\ldots,T_n} & = & T_i
\end{array}
\]
    \caption{Substitution of type variables}
    \label{fig:substitution}
\end{figure}

Substitution replaces every occurrence of a type variable, \emph{including} the ones in the method signatures of a method-local class.
Substitution does not rename the type variables it meets.
No renaming is needed: the substituted types are ground, and the domain of the substitution never contains a type variable bound inside the substituted term, by requirement A7 in the program and by $\mathsf{fresh}$ during reduction.
In the domain of a substitution, bounds $\Delta$ stand for $\mathrm{dom}(\Delta)$, and funnels $\pl{F}$ stand for their funnelled type variables:
$\tsubst{T}{F_1,\ldots,F_n}{T_1,\ldots,T_n}$ replaces $X_i$ with $T_i$ for every $F_i=\funnel{X_i}$, and ignores every $F_i=G_i$.
Substitution of type variables into expressions happens only during reduction, where the substituted types are ground.
Thus $\tsubst{\funnel{X_i}}{X_1,\ldots,X_n}{T_1,\ldots,T_n}=T_i$ is a ground type $G$, and the result is again a funnel.

\begin{figure}[H]\footnotesize
    \centering
    \setlength{\arraycolsep}{3pt}
\[
\newcommand{\s}[1]{\tsubst{#1}{\pl{x}}{\pl{v}}}
\begin{array}{rcl}
\s{x} & = & x \quad\textrm{if }x\not\in\pl{x}\\
\tsubst{x_i}{x_1,\ldots,x_n}{v_1,\ldots,v_n} & = & v_i\\[1ex]
\s{\call{e}{\pl{T}}{m}{\pl{e}'}} & = & \call{\s{e}}{\pl{T}}{m}{\s{\pl{e}'}}\\[1ex]
\s{\explicitcall{e}{T}{\pl{T}}{m}{\pl{e}'}} & = & \explicitcall{\s{e}}{T}{\pl{T}}{m}{\s{\pl{e}'}}\\[1ex]
\s{e\fstop f} & = & \s{e}\fstop f\\[1ex]
\s{\obj{C}{\pl{T}}{\pl{e}}} & = & \obj{C}{\pl{T}}{\s{\pl{e}}}\\[1ex]
\twolines{\kw{if}\;\OR e_0\;\kw{instanceof}\;T\;x'\CR}{\OC\;e_1\;\CC\;\kw{else}\;\OC\;e_2\;\CC\;\s{}} & = &
  \twolines{\kw{if}\;\OR\s{e_0}\;\kw{instanceof}\;T\;x'\CR}{\OC\;\s{e_1}\;\CC\;\kw{else}\;\OC\;\s{e_2}\;\CC}\\[3ex]
\twolines{\kw{new}\;\gen{C}{\pl{F}}\;\kw{extends}\;\gen{C'}{\pl{T}}}{\OR\pl{e}\CR\OC\pl{M}\CC\;\s{}} & = &
  \twolines{\kw{new}\;\gen{C}{\pl{F}}\;\kw{extends}\;\gen{C'}{\pl{T}}}{\OR\s{\pl{e}}\CR\OC\s{\pl{M}}\CC}\\
  && \quad\textrm{if }\this\not\in\pl{x}\\[1ex]
\twolines{\kw{new}\;\gen{C}{\pl{F}}\;\kw{extends}\;\gen{C'}{\pl{T}}}{\OR\pl{e}\CR\OC\pl{M}\CC\;\tsubst{}{\pl{x},\this,\pl{x}'}{\pl{v},v,\pl{v}'}} & = &
  \twolines{\kw{new}\;\gen{C}{\pl{F}}\;\kw{extends}\;\gen{C'}{\pl{T}}}{\OR\tsubst{\pl{e}}{\pl{x},\this,\pl{x}'}{\pl{v},v,\pl{v}'}\CR\OC\tsubst{\pl{M}}{\pl{x},\pl{x}'}{\pl{v},\pl{v}'}\CC}\\[3ex]
\twolines{\methodsig{\Delta}{T_0}{m}{T_1\;x'_1,\ldots,T_k\;x'_k}}{\OC\;\ret{e}\;\CC\;\s{}} & = &
  \twolines{\methodsig{\Delta}{T_0}{m}{T_1\;x'_1,\ldots,T_k\;x'_k}}{\OC\;\ret{\s{e}}\;\CC}
\end{array}
\]
    \caption{Substitution of variables}
    \label{fig:termsubstitution}
\end{figure}

A method-local class binds \this, as in Java: inside the methods of a method-local class, \this is the local object, and the outer \this is hidden.
Thus substitution of variables enters the methods of a method-local class for every variable except \this.
The variable $x'$ of an \kw{if} and the parameters $x'_1,\ldots,x'_k$ of a method are never in $\pl{x}$, by requirement A7 in the program and by $\mathsf{fresh}$ during reduction; so no renaming is needed.

\ldots

The rule \textsc{new-local} is the heart of funnelling, including subtle management of type environments.
The rule constructs a restricted set of type bounds and variable bindings reflecting only the explicitly-funnelled type variables:
$\Delta'$ includes only the explicitly-funnelled type variables, and $\Gamma'$ keeps only the variables of $\Gamma$, other than \this, whose types mention only type variables of the result type $\gen{C}{|\pl{F}|}$ (which must be well-kinded under $\Delta'$).
The bounds $\Delta'$ are computed by $\mathsf{scope}(\Delta;\pl{F};|\pl{F}|)$, which is undefined when a funnelled type variable is not in $\Delta$.
However, the rule still uses the \emph{full} ambient environments for checking certain subexpressions. This is because some $e\in\pl{e}$ may initially contain uses of \emph{non}-funnelled type variables (e.g., an explicit call), but will reduce to a value which no longer depends on such type variables.
This is why the premise $\typing{\Delta}{\Gamma}{\obj{C'}{\pl{T}}{\pl{e}}}{\_}$ checks the effective supertype instance under $\Delta;\Gamma$, but the premise $\kind{\Delta'}{\gen{C'}{\pl{T}}}$ ensures the resulting type of that
expression satisfies the funnelling constraints by checking kinding under the more restrictive $\Delta'$.

\begin{figure}[H]\small
    \centering
    \begin{mathpar}
\inferrule[new-local-java]{
  \typing{\Delta}{\Gamma}{\obj{C'}{\pl{T}}{\pl{e}}}{\_}\\
  \kind{\Delta}{\gen{C'}{\pl{T}}}\\\\
  %\kind{\Delta}{C\langle\rangle}\\
  \mok{\Delta}{\Gamma\setminus\this,\gen{C}{}\;\this}{\pl{M}}\\\\
  \forall C'',\pl{T}'\ldotp\subtype{\Delta}{\gen{C'}{\pl{T}}}{\gen{C''}{\pl{T}'}}\Rightarrow
    \mathsf{overrideOk}(\gen{C''}{\pl{T}'},\pl{M})
}{
  \typing{\Delta}{\Gamma}{
    \innerobj{C}{}{C'}{\pl{T}}{\pl{e}}{\pl{M}}
  }{\gen{C}{}}
}
    \end{mathpar}
    \caption{Java typing of local classes}
    \label{fig:java}
\end{figure}

Rule \textsc{new-local-java} of Figure~\ref{fig:java} types local classes as Java does.
A local class has no funnels, and its methods are typed under the full $\Delta$ and $\Gamma$ in scope, except for \this.
Replacing rule \textsc{new-local} with rule \textsc{new-local-java} makes our system unsound for the same reason Java is unsound.
Consider the method \QQ{mk} below, the method \QQ{m} of Example~1 in our syntax:
\begin{lstlisting}
<X> Foo<X> mk(Object<> other, X x){ return
  if (other instanceof MyFoo<> w){ w }
  else { new MyFoo<> extends Foo<X>(){ X x(){ return x; } } }; }
\end{lstlisting}
Rule \textsc{new-local-java} gives to the object created by \QQ{mk} the type \QQ{MyFoo<>}.
Rule \textsc{sub-inherit-local} reads the declaration of \QQ{MyFoo}, so \QQ{MyFoo<>} is a subtype of \QQ{Foo<X>}, where \QQ{X} is the type variable of \QQ{mk}.
The type \QQ{MyFoo<>} does not say which type \QQ{X} stood for when the object was created.
Since $\generic{\QQ{MyFoo}}=\QQ{MyFoo<>}$ has no type variables, rule \textsc{cast-down} accepts the type test \QQ{other instanceof MyFoo<>} from \QQ{Object<>}.
Thus the call \QQ{a.<String>mk(a.<Integer>mk(o, i), s)} has type \QQ{Foo<String>}, but it returns a \QQ{Foo<Integer>}.
This is the unsoundness of Java.
It is there even if we keep the restriction of Java that the name of a local class can only be written inside the method declaring it, as \QQ{mk} does.
Our formal system shows two more failure points.
They are the same failure, shown more explicitly:
\begin{itemize}
\item Reduction does not preserve types.
The call \QQ{a.<Integer>mk(o, i)} has type \QQ{Foo<Integer>}.
The value \QQ{o} is not a \QQ{MyFoo}, so the type test fails and the call takes the \kw{else} branch.
It reduces to the value
\begin{lstlisting}
new MyFoo<> extends Foo<Integer>(){ Integer x(){ return i; } }
\end{lstlisting}
The type of this value is \QQ{MyFoo<>}, and \QQ{MyFoo<>} is not a subtype of \QQ{Foo<Integer>}.
The supertypes of \QQ{MyFoo<>} come from the declaration of \QQ{MyFoo} in \QQ{mk}, not from the value:
rule \textsc{sub-inherit-local} derives $\subtype{\emptyset}{\QQ{MyFoo<>}}{\QQ{Foo<X>}}$, where \QQ{X} is the type variable of \QQ{mk}.
Subtyping does not check kinding, so this judgement holds in the empty bounds $\emptyset$, where \QQ{X} is not in scope.
However, $\QQ{Foo<X>}$ is not well kinded under $\emptyset$, so rule \textsc{subs} cannot give the value the type \QQ{Foo<X>} either.
Java has no typing of values during reduction, so in Java this failure is invisible.
\item Typing depends on the names of type variables.
Our system lets the name of a local class be written anywhere.
Consider the methods
\begin{lstlisting}
<Y> MyFoo2<> mk2(Y y){
  return new MyFoo2<> extends Foo<Y>(){ Y x(){ return y; } }; }
<Y> Foo<Y> leak(MyFoo2<> z){ return z; }
\end{lstlisting}
The method \QQ{leak} is well typed, since the type variable \QQ{Y} of \QQ{leak} has the same name as the type variable \QQ{Y} of \QQ{mk2}.
Renaming \QQ{Y} into \QQ{Z} in \QQ{leak} makes \QQ{leak} ill typed.
The call
\begin{lstlisting}
a.<String>leak(a.<Integer>mk2(i))
\end{lstlisting}
has type \QQ{Foo<String>}, but it returns a \QQ{Foo<Integer>}, without type tests.
\end{itemize}
Funnelling solves the problem in all three forms.
Under rule \textsc{new-local}, the object created by \QQ{mk} has type \QQ{MyFoo<X>}, written \QQ{new MyFoo<X funnel> extends Foo<X>()\{..\}}.
Since $\generic{\QQ{MyFoo}}=\QQ{MyFoo<X>}$, rule \textsc{cast-down} accepts the type test \QQ{instanceof MyFoo<X>} from \QQ{Foo<X>}, but not from \QQ{Object<>}.
Thus the funnelled \QQ{mk} takes \QQ{Foo<X> other} instead of \QQ{Object<> other}, and tests \QQ{other instanceof MyFoo<X> w}.
The call \QQ{a.<Integer>mk(o, i)}, with \QQ{o} of type \QQ{Foo<Integer>}, reduces to a value of type \QQ{MyFoo<Integer>}, a subtype of \QQ{Foo<Integer>}.

\ldots

\todo[inline]{briefly note that Java allows extra, ``fresh'' generic parameters in local classes}

\todo[inline]{call out that we're allowed to declare return types that mention an inner class; inner classes are ``public'', and our approach works even in this case, the more restrictive case of Java is naturally covered}

\todo[inline]{articulate early:
    It's not just about generic reification: just checking type parameters aren't enough, as our MyFoo cases have no generic params.
    You need to *deeply* check reified generics, you need to check not just the base type but also check the choice(s) of parent types from the context into the runtime check.
}

\begin{figure}[H]
\centering
\[
\begin{array}{rrcl}
\textrm{Values} & v & ::= & \obj{C}{\pl{G}}{\pl{v}} \mid \innerobj{C}{\pl{G}}{C'}{\pl{G}'}{\pl{v}}{\pl{M}}\\
\textrm{Eval. Ctx.} & \mathcal{E} & ::= & \square \mid \call{\mathcal{E}}{\pl{G}}{m}{\pl{e}} \mid \call{v}{\pl{G}}{m}{\pl{v},\mathcal{E},\pl{e}}\\
    &&& \mid \explicitcall{\mathcal{E}}{G}{\pl{G}}{m}{\pl{e}} \mid \explicitcall{v}{G}{\pl{G}}{m}{\pl{v},\mathcal{E},\pl{e}}\\
    &&& \mid \mathcal{E}\fstop f \mid \obj{C}{\pl{G}}{\pl{v},\mathcal{E},\pl{e}}\\
    &&& \mid \innerobj{C}{\pl{G}}{C'}{\pl{G}'}{\pl{v},\mathcal{E},\pl{e}}{\pl{M}}\\
    &&& \mid \ifthen{\mathcal{E}}{G}{x}{e_1}{e_2}\\
\textrm{Contexts} & \mathcal{C} & ::= & \square \mid \call{\mathcal{C}}{\pl{T}}{m}{\pl{e}} \mid \call{e}{\pl{T}}{m}{\pl{e},\mathcal{C},\pl{e}'}\\
    &&& \mid \explicitcall{\mathcal{C}}{T}{\pl{T}}{m}{\pl{e}} \mid \explicitcall{e}{T}{\pl{T}}{m}{\pl{e},\mathcal{C},\pl{e}'}\\
    &&& \mid \mathcal{C}\fstop f \mid \obj{C}{\pl{T}}{\pl{e},\mathcal{C},\pl{e}'}\\
    &&& \mid \innerobj{C}{\pl{F}}{C'}{\pl{T}}{\pl{e},\mathcal{C},\pl{e}'}{\pl{M}}\\
    &&& \mid \ifthen{\mathcal{C}}{T}{x}{e_1}{e_2}\\
    &&& \mid \ifthen{e_0}{T}{x}{\mathcal{C}}{e_2}\\
    &&& \mid \ifthen{e_0}{T}{x}{e_1}{\mathcal{C}}
\end{array}
\]
\caption{Values and contexts}
\label{fig:contexts}
\end{figure}

Marco had this idea that casts may be avoided by using dynamic dispatch returning this.
Can not work, since Meth local class is private, there is no way an external interface can mention it to provide the cast method.
Of course, the method local is private to the body (this means also the method signature can not use it)
This means no self:Foo method, nor self:T method of a generic class eventually instantiating T=Foo.

This means the core language DO NEED if(e instanceof T x){} else {}

//nested funnelling:
  new A<X>(){ B<X> f()-> new B<X>(){} }
  new A<X>(){ <K>B<X,K> f()-> new B<X,K>(){} }

using 'extends' for declarations in top level classes and generic methods
using 'funnel' for named inner classes
  new A<X funnel Object>(){ B<X> f()-> new B<X funnel Object>(){} }
  new A<X funnel Object>(){ <K extends Object> B<X,K> f()-> new B<X funnel Object,K funnel Object>(){} }
  or, if the bounds are not repeated
  new A<X funnel>(){ <K extends Object> B<X,K> f()-> new B<X funnel,K funnel>(){} }
Very funky syntax:
we get new B<X funnel Object,K funnel Object>(){}
but, during reduction, this can become 
we get new B<Person,K funnel Object>(){}
where we have a mixture of concrete types and generic funnelling

Global class definition uses class C<X extends Bar>{..}
named method local class uses new C<X funnel>(){..}
normal new: new C<Ts>()

\begin{verbatim}
e ::= x | e.m(es) | e.f | new C<Ts>(es) | new C<Fs>(){Ms} | if (e0 instanceof T x){ e1 }else{ e2 }
T ::= C<Ts> | X
F ::= C<CCs> | X funnel// Important, can not be T | X funnel, otherwise we could have an X not 'funnelled'
CC::= C<CCs> //types with zero free type variables
B::=  X extends T // Need T, not F, to be able to do things like (Java) class Bar<X,Y extends X> {}
Dec::= class C<Bs>{ Ms Ss} //assume standard constructor
S ::= T x; //field/slot
M ::= Sig; | Sig{ return e; }
Sig::= <Bs> T m(x1:T1,..xn:Tn)
G ::= x1:T1..xn:Tn
\end{verbatim}
//we do not need Object to be able to write 'extends Object'
//this is exactly as the distinction of checked and unchecked exceptions: formally it evaporates in syntacitc sugar where all is checked
//and all methods 'throws Error, RTE'
// Can assume the user provides their own Object type
Note: an use of X is just X, a declaration of X is X extends T; a funnelling of X is X funnel
Program is just many Dec plus an e typable in empty
MARCO TODO: 
  typechecking
    Decs;Bs;G|-e:T
    Decs|-Dec:ok
    Decs;T|-M:ok
    e --> e'
Note, there are multiple semantics/ts:
  e instanceof CC //always clear
  e instanceof T with T not of form CC would not typecheck

  e --> e' for 
  Java with warn->err: just have CC in instanceof
  Java with magic reifined cast: just have T in instanceof (since reduction will replace(thus remember) generics
  Java with warn->off:
        AT RUN TIME instanceof only checks base class, ignores generic parameters
        AT COMPLILE TIME, instanceof can be err if the full type (including generics) is not potentially correlated)?? unobvious consider interfaces
        In prose, should explain that we won't flag cast of some type Foo to some type Bar, even though Java sometimes flags unrelated casts
\mjnote{these three semantics are replaced by the single one of Java 16+: erased run-time check, plus the checked-cast premise of rule \textsc{if-instanceof} (Figure~\ref{fig:cast}). Unrelated types are rejected by it, as in Java.}

<X extends List<Y funnel>>

class Foo<X> {
    class Getter<X funnel> {
        <Y extends List<X>> X getFirst(Y y) { y.getOrElse(0,null) }
    }
}

<Y> .. return foo<Y,MyX>();
<Y,X extends List<Y>> foo(){ return Foo<X funnel>(){... X }; }
